% Fuente: Auxiliar 6, Pregunta 3.
La matriz A es:
\[
\begin{pmatrix}
    1 & -2 & 5 & 2 & 0 & -3 & 0 \\
    0 & 1 & -2 & 4 & 1 & -2 & 0 \\
    0 & 1 & 0 & -3 & 0 & -3 & 1
\end{pmatrix},
\]

\textbf{I iteración:}
\begin{enumerate}
    \item Se incia con una base formada por $(A_1|\: A_5 | \: A_7)$ y $x_B'=(x_1 \quad x_5 \quad x_7)$.

\[
B^{-1}=
\begin{pmatrix}
    1 & 0 & 0 \\ 0 & 1 & 0 \\ 0 & 0 & 1
\end{pmatrix}, \quad
x_B=B^{-1}b=\begin{pmatrix}
    5 \\ 7 \\ 3
\end{pmatrix}
\]

\item  Se calculan los costos reducidos $\overline{c_j} =c_j - c_B B^{-1} A_j$. Notar que $c_B'=(0\quad 0 \quad 0)$
\[
\begin{aligned}
    &\overline{c_2}=-3\\
    &\overline{c_3}=-2 \\
    &\overline{c_4}=1 \\
    &\overline{c_6}=0
\end{aligned}
\]
Luego $x_2$ y $x_3$ pueden entrar a la base, ya que mejoran la función objetivo. Se selecciona $x_2$, para entrar a la base.
\item Cálculo de $u$.
\[
u=B^{-1} A_2=\mathbb{I} \begin{pmatrix}
    -2 \\ 1 \\ 1
\end{pmatrix}
\]
Dado que las componentes $u_2$ y $u_3$ son positivas, $x_5$ y $x_7$ son candidatos a salir de la base.

\item Cálculo de $\theta^*$, para ver la variable que saldrá:
\[
\theta^*=\min\left( \frac{x_5}{u_5}, \frac{x_7}{u_7}\right)= \min(7,3)=3
\]
luego $x_7$ saldrá de la base; como se ubica en tercera posición de la base, se tiene que $l=3$.

\item  Se calcula el valor de $\overline{x_B}$ para la siguiente iteración.
\[
B^+=\{x_1,x_5,x_2\},
\qquad
x_B^+
=
\begin{pmatrix}
x_1\\
x_5\\
x_2
\end{pmatrix}
=
\begin{pmatrix}
11\\
4\\
3
\end{pmatrix}.
\]

\item Se obtiene $\overline{B^{-1}}$ para la siguiente iteración. Para ello se forma $(B^{-1}|u)$ y se debe pivotar utilizando la fila $l=3$.
    \[
\left(\begin{array}{ccc|c}
1 & 0 & 0 & -2 \\
0 & 1 & 0 & 1 \\
0 & 0 & 1 & 1
\end{array}\right)
\;\rightarrow \;
\left(\begin{array}{ccc|c}
1 & 0 & 2 & 0 \\
0 & 1 & -1 & 0 \\
0 & 0 & 1 & 1
\end{array}\right) \; \rightarrow \; \overline{B^{-1}} = \begin{pmatrix}
    1 & 0 & 2 \\ 0 & 1 & -1 \\ 0 & 0 & 1
\end{pmatrix}
\]
Con los valores de $\overline{B^{-1}}$ y $\overline{x_B}$, se termina la I iteración y se comienza la siguiente.
\end{enumerate}

En relación con la segunda pregunta, se tiene que la función objetivo:
\begin{itemize}
    \item Mejora para $x_2$ y $x_3$
    \item Empeora para $x_4$
    \item No cambia para $x_6$
\end{itemize}
La función objetivo mejorará al agregar a la base una variable no básica cuyos costos reducidos sean $x_j<0$. Los costos reducidos muestran la tasa de cambio de la función objetivo al moverse de $x$ a $x-\theta^*u$.

Los costos reducidos muestran la tasa de cambio de la función objetivo al moverse desde la solución básica actual en la dirección asociada a la variable entrante.
